\chapter{The New Soviet Order}

The coming of NEP, which had the unpremeditated consequence of strengthening the central authority of the party, also encouraged the centralizing forces already at work in the formation of the Soviet state. The mass enthusiasm of 1917 for the destruction of state power had faded into the world of unrealized dreams. These continued to haunt the memories of many party members. But since Brest-Litovsk, and since the civil war, the need to create a state power strong enough to cope with these emergencies had been perforce accepted; and it was now reinforced by the need to rebuild the devastated and shattered national economy. The NEP period not only shaped what was to become the permanent constitutional structure of the USSR, but determined the lines to be followed for many years in its relations with foreign countries. 
% 注：去除了 "land since", "Istrong", "(accepted" 中的边缘乱码符号。

It was time to stabilize the fluid constitutional arrangements of the Soviet regime. A constitution of the Russian Socialist Federal Soviet Republic (RSFSR) had been promulgated in July 1918. It opened with the ``Declaration of Rights of the Toiling and Exploited People'', proclaimed by the All-Russian Congress of Soviets six months earlier (see p. 9 above). It vested supreme authority in an All-Russian Congress of Soviets composed of delegates elected by city and provincial Soviets, representation being heavily weighted in favour of the cities, the home of the workers. The franchise was confined to those who ``earn their living by production or socially useful labour'', together with soldiers and disabled persons. The congress elected an All-Russian Central Executive Committee (VTsIK) to exercise authority on its behalf in the intervals between sessions; and VTsIK in turn appointed a Council of People's Commissars (Sovnarkom), whose main functions were administrative, but which was also entitled to issue orders and decrees, so that no clear line of demarcation was drawn between the powers of Sovnarkom and VTsIK. The constitution also enunciated such general principles as the separation of church and state; freedom of speech, opinion and assembly for the workers; the obligation for all citizens to work on the principle, ``He who does not work, neither shall he eat''; liability to military service for the defence of the republic; and abolition of all discrimination on grounds of race or nationality. The chaos of the civil war precluded any definition of the territory of the republic. The term ``federal'' in the title of the republic had no exact meaning; it covered both the incorporation in the RSFSR of ``autonomous'' republics and regions inhabited mainly by non-Russian populations, and the establishment of links between the RSFSR and other Soviet republics which had been, or would be, proclaimed in other parts of the former Russian Empire. These links took at first the form of alliance rather than federation. The RSFSR concluded treaties of alliance with Azerbaijani and Ukrainian Soviet republics in September and December 1920, and with Belorussian, Armenian and Georgian republics in 1921. Resistance to the process of unification was encountered in the Ukraine, where a national anti-Soviet government had been one of several rival authorities during the civil war, and in Georgia, where a Menshevik government had installed itself. Military power was used to expel the dissidents, and set up unimpeachably Bolshevik governments. The use of force could be more easily justified in the Ukraine, which had been deeply involved in the civil war, and where rival armies had reduced much of the country to anarchy, than in Georgia, which long remained a restive and recalcitrant member of the federation of Soviet republics. 
% 注：去除了 "^Congress", "/in favour", "nd assembly" (补全为 and), "f /December", "(rival" 中的杂乱标记。

As the whole country advanced towards economic recovery, and sought to renew contacts with the outside world, it seemed natural and necessary that it should function for these purposes as a single unit. While the form, and some of the reality, of local autonomy was carefully pursued, the Russian Communist Party, to which regional parties were affiliated, maintained a uniform discipline; and major decisions of economic and foreign policy were taken in Moscow. The first step was to persuade the three Transcaucasian republics---Armenia, Georgia and Azerbaijan---to unite in a Transcaucasian Socialist Federal Republic. Then, in December 1922, congresses of the four republics---the RSFSR and the Ukrainian, Belorussian and Transcaucasian republics---meeting separately, voted to form a union of Soviet socialist republics (USSR). Finally, delegates of the four republics met together, constituted themselves as the first Congress of Soviets of the USSR, and elected a committee with instructions to draft a constitution. The constitution of the USSR was approved by the committee in July 1923, and formally ratified by the second Congress of Soviets of the USSR in January 1924. 
% 注：去除了 "land sought", "TTTt Tftt", "I seemed", "f these" 等干扰符号。

The constitution was modelled on the original constitution of the RSFSR. The sovereign Union Congress of Soviets was composed of delegates of the congresses of Soviets of the constituent republics, representation being proportionate to the population of each republic. The congress elected a central executive committee (TsIK), which appointed a Sovnarkom of the USSR. The organization of People's Commissariats was complicated. Foreign affairs, foreign trade, military affairs, and ``the conduct of the struggle against counter-revolution'' by the Cheka, now renamed OGPU (``Unified State Political Administration''), were reserved exclusively to the Union authorities; each republic had its GPU, which was, however, directly subordinate to the OGPU. Most economic affairs were subject to a system of ``unified'' commissariats; there were commissariats of the Union and commissariats of the republics, the latter enjoying a certain degree of independence. In other fields of administration, including agriculture, internal affairs, health and education, the republics alone had commissariats without any Union counterpart. In form the USSR was a federation of republics. But the omission of the word ``federal'' from its title was perhaps significant, since its unifying tendencies were apparent from the outset. The RSFSR accounted for more than 75 per cent of the population of the Union, and 90 per cent of its area. The other republics had some reason to suspect that the USSR was little more than the RSFSR writ large, and represented an extension over them of the central authority of Moscow. Voices of dissent were heard, especially from Ukrainian and Belorussian delegates, in the committee which drafted the new constitution. 
% 注：去除了 "II affairs", "Isariats", "'was a", "nad some", "/than the", "/over them" 等杂乱的识别符。

An attempt to meet these objections led to one notable innovation, designed to recognize the formal equality of all the republics. The TsIK of the USSR was divided into two chambers. The first and much larger chamber, the Council of the Union, consisted of delegates elected proportionately to the population of the republics; this recognized the enormous preponderance of the RSFSR. Delegates to the second chamber, the Council of Nationalities, were elected on the basis of equality of national groups, five from each of the four main republics and each of the autonomous republics, one from each autonomous region. But, since both chambers normally met only to listen to, and comment favourably on, statements of official policy (on occasion they met jointly to hear important speeches), and since contentious issues were seldom raised, and never voted on, these complicated arrangements had no practical significance in the process of policy-making. The periodical sessions of the congress and of TsIK, whose membership was enlarged as time went on, took no decisions. But they provided an important means of making contact with representatives of the outlying and often primitive regions of the Union, and of popularizing and making known throughout the Union major policies decided in Moscow; their principal function was not to debate, but to instruct, persuade and exhort. The constitutions of the USSR and of its constituent units served quite different purposes from the constitutions of western countries, to which they had only the most superficial resemblance. 
% 注：清理了 "-lUU III", "[ jseldom", ".contact", "'their" 等无意义的标点和字母。

The complex ethnic structure of Central Asia, and its affiliations with the Muslim world, made it a special problem. The Central Asian republics of Bokhara and Khorezm, though they had been brought within the orbit of Moscow by treaties of alliance with the RSFSR, were excluded from these constitutional arrangements on the ground that they were not yet socialist. It was not till 1925 that Central Asia was re-organized on national lines; Uzbek and Turkmen Socialist Soviet Republics, with subordinate autonomous units, were incorporated in the USSR as its fifth and sixth constituent republics. 

The structure of the party organization was no less important a factor in the course of events than the structure of the Soviets. Between party congresses the supreme authority was vested in the party central committee. The committee which took the vital decisions to launch the rising in October 1917, and later to sign the Brest-Litovsk treaty, consisted of 22 members. In the period of acute crisis which followed, this body was too unwieldy for rapid action, and decisions on crucial issues rested in practice with Lenin in consultation with other top leaders. The eighth party congress in March 1919 elected a central committee of nineteen full members with eight candidates, who were entitled to attend meetings, but not to vote. But it appointed a Politburo of five, to be responsible for political decisions, and an Orgburo to take control of questions of party organization; and this marked the atrophy of the central committee as an effective source of authority. The ninth congress in 1920 reorganized the secretariat, placing it under the management of three ``permanent'' secretaries, who were members of the party central committee; and in the ensuing period the refurbished secretariat underwent a rapid expansion, acquiring a staff of several hundred officials, divided into departments charged with different branches of party activity. The party structure assumed the shape which it was to retain throughout the nineteen-twenties, though the processes at work took some years to develop to their full extent. The creation of a powerful party machine later provided an instrument for Stalin's dictatorship. Party congresses met annually till 1925, and thereafter at less regular intervals, alternating with smaller and less formal party conferences; and the party central committee held three or four sessions a year. These bodies continued to serve as a forum where important issues were debated, though manipulation by the secretariat of the election of delegates soon made the results a foregone conclusion. Only the Politburo, increased in number to seven, and later to nine, with several candidates, remained throughout the nineteen-twenties a source of decisions at the highest level; and, since the authority of the party in a one-party state was mandatory for all decisions and activities of the Soviet Government, the party Politburo became the supreme policy-making organ of the USSR. 
% 注：去除了 "^iBoviets", "jlvested" 以及断字导致的语序错位 ("I be responsible... 1919 elected... but not to vote. But it appointed...")，并恢复正常逻辑。

The strengthening of party and Soviet organization was matched by a consolidation of Soviet relations with the outside world. Even in the days of war communism, when thoughts of world revolution were uppermost in Moscow, the rare opportunities of direct contact with the governments of western countries were not neglected. In January 1920 representatives of the Russian cooperatives in Paris discussed with representatives of Western governments the resumption of an exchange of goods with Soviet Russia; and Litvinov in Copenhagen negotiated an agreement for the mutual repatriation of prisoners. A treaty of peace with Estonia was signed on February 12, 1920; and Lenin commented that ``we have opened a window on Europe, which we shall try to utilise as extensively as possible''. At the party congress in March 1920 Lenin spoke of the need ``to manoeuvre in our international policy''. A few days later Krasin, the one leading Bolshevik who had foreign industrial and commercial experience, set out with a large delegation of ``trade experts'' for Scandinavia, and in May was politely received in London. These overtures were cut short by the Polish war, which inspired a recrudescence of revolutionary hopes in Moscow and a fresh outbreak of apprehension and animosity in the west. But by the autumn of 1920 peace was restored. A Russian trading company was registered in London under the name Arcos (All-Russian Cooperative Society); and Krasin spent most of the winter in London negotiating with the British Government and with firms interested in orders for Soviet Russia. Finally, just a week after Lenin had introduced NEP to the party congress in Moscow, an Anglo-Soviet trade agreement was signed in London on March 16, 1921. 
% 注：修复了由于分栏原因造成的句子交织 ("12, 1920; and Lenin commented that “we have opened a negotiated an agreement for the mutual repatriation of prisoners. A treaty of peace with Estonia was signed on February ^ndow on Europe, which we shall try to utilise as extensively 'as possible”") 并删除了乱码。

The trade agreement was rightly hailed as a break-through and a turning-point in Soviet policy. The parties agreed to put no obstacles in the way of trade with one another, and, in default of formal diplomatic recognition, to exchange official trade representatives. The most important clause from the British point of view was one in which each party undertook to ``refrain from action or undertakings'' and from ``official propaganda, direct or indirect'', against the other. ``Action or propaganda to encourage any of the peoples of Asia in any form of hostile action against British interests or the British Empire'' was specially mentioned. A pledge to refrain from hostile propaganda had been given in a less elaborate form in the Brest-Litovsk treaty. But the circumstances were different. That treaty had been concluded in conditions which were not expected to last, and did not last. The Anglo-Soviet agreement was designed, like NEP, ``seriously and for a long time''. It heralded a change of emphasis in Soviet policy. Pronouncements about world revolution continued to be made, but were consciously or unconsciously regarded more and more as a prescribed ritual, which did not affect the normal conduct of affairs. The latent incompatibility between the policies of the People's Commissariat of Foreign Affairs (Narkomindel) and of Comintern began to come to the surface. 

The background of the Soviet rapprochement with Great Britain was economic; desire to facilitate mutually profitable trade. The background of the rapprochement with Germany was primarily political, being rooted in common opposition to the Versailles treaty and common antipathy to the claims of Poland. Radek, who spent most of 1919 in prison or under house-arrest in Berlin, contrived to make contact with Germans from many different milieux, to all of whom he preached the virtues of German-Soviet cooperation. Official German-Soviet relations had been severed since the assassination of the German Ambassador in Moscow in 1918. In the summer of 1920 a Soviet representative was once more received in Berlin, and a German representative in Moscow. The Polish war gave a powerful stimulus to friendly relations between Poland's two neighbours. Trotsky was reported to favour an agreement with Germany; and Lenin in a public speech in November 1920 noted that, ``though the bourgeois German Government madly hates the Bolsheviks'', nevertheless ``the interests of the international situation are pushing it towards peace with Soviet Russia against its will''. Soviet policy was still ambivalent, being divided between the pursuit of revolution and diplomacy. On March 17, 1921, the German Communist Party began an armed rising against the government, known in party history as the ``March action''. The enterprise was certainly supported, perhaps prompted, by Zinoviev and Comintern officials; the involvement of the other Soviet leaders, heavily preoccupied at the moment by the Kronstadt revolt and the party congress, is doubtful. But the defeat of the German rising must have further diminished the waning hopes in Moscow of revolutions in the west, and strengthened the hand of those who saw a diplomatic accommodation with the capitalist countries as the immediate goal. 
% 注：去除了 "■'Radek", "!!with Germany", "[920 noted", "Hlent", "/diplomacy", "[in Moscow", "ffiand of those" 等边缘标记。修复了 "November [920" 为 "November 1920"。

A feature of German-Soviet relations at this time was the quest for military collaboration, arising out of the prohibition under the Versailles treaty on the manufacture of armaments in Germany. In April 1921 Kopp, the Soviet representative in Berlin, after secret discussions with the Reichswehr, brought back to Moscow a scheme for the manufacture in Soviet Russia by German firms of guns and shells, aeroplanes and submarines. The response was favourable, and a German military delegation visited Moscow during the summer. An agreement was clinched at meetings in Berlin in September 1921 at which Krasin and Seeckt, the head of the Reichswehr, were the principal negotiators; it seems to have been at this moment that Seeckt first divulged to the civilian German Government what was on foot. The submarine project was dropped. But German factories in Soviet Russia were soon engaged in the production of guns, shells and aeroplanes. Tanks were added to the programme; and experiments were conducted in gas warfare. The products of these enterprises were supplied both to the Reichswehr and to the Red Army. Later German officers trained Red Army personnel in tank warfare and in military aviation. These arrangements were veiled in the utmost secrecy. No mention was made of them in the Soviet press; and they were for a long time successfully concealed from the German public and German politicians, as well as from the western Allies. It was a far cry from the days when, on the morrow of the revolution, the Bolsheviks had denounced the secret treaties concluded by the Tsarist government with the Allies during the war. Meanwhile, German-Soviet economic relations were cemented by the setting up of ``mixed companies'' and the granting of ``concessions'' in Soviet Russia to German firms. 
% 注：去除了 "Iman-Soviet" 中的识别符。

Early in 1922 both the Soviet and the German Governments were invited to attend an international conference, which met at Genoa on April 10. The conference was a bold attempt by Lloyd George, its most active promoter, to re-forge links with Germany and Soviet Russia, hitherto outcasts from the European community. Lenin greeted the invitation with guarded enthusiasm. ``We go to it'', he explained, ``as merchants, because trade with the capitalist countries (so long as they have not altogether collapsed) is unconditionally necessary for us, and we go there to discuss . . . the appropriate political conditions for this trade''. Chicherin, Krasin and Litvinov led the Soviet delegation, the first of its kind to attend an international conference on equal terms with the delegations of other major Powers. The conference was a failure, partly owing to unyielding French opposition to Lloyd George's aims, partly owing to the inability of British and Soviet negotiators to reach agreement on the question of Soviet debts and liabilities. The Soviet Government was ready in principle to recognize pre-war debts (though not the war debts) of the former Russian Government, but only provided a substantial foreign loan were granted to facilitate their settlement. The Soviet Government refused to rescind decrees nationalizing foreign enterprises, but was prepared in certain conditions to allow foreign firms to re-occupy their former enterprises in the guise of ``concessions''. No amount of ingenuity could bridge these gaps. 

The deadlock in negotiations paradoxically produced the only concrete result of the conference. For some time, Soviet and German diplomats in Berlin had been discussing the terms of a political treaty. The Soviet delegation at Genoa, having failed to make any impression on the western Allies, now pressed the German delegation, headed by the Foreign Minister, Rathenau, to complete and sign the treaty forthwith; and the German delegation, equally disillusioned by the proceedings of the conference, agreed. The treaty was signed, hastily and secretly, at Rapallo on April 16, 1922. The contents of the Rapallo treaty were not remarkable. The only operative clauses provided for a mutual renunciation of financial claims, and the establishment of diplomatic and consular relations. But, as a demonstration of solidarity against the Western Allies, it shattered the conference, and made a lasting impact on the international scene. Soviet Russia had secured for herself a bargaining position among the European Powers. Manoeuvres originally conceived as expedients to tide over a crisis were becoming accepted procedure. 
% 注：清除了 ".Minister", "'ceedings", ".scene" 等多余的点和符号。

In Comintern, signs of the change of mood were apparent as early as its third congress in June 1921. The effervescent revolutionary enthusiasm of the second congress a year earlier had evaporated. What the Bolsheviks initially thought of as impossible had happened: the socialist Soviet republic had maintained itself, and showed every sign of continuing to maintain itself, in a capitalist environment. Lenin found himself on the defensive at the congress both in domestic and in international affairs. He took pains to explain the necessity of NEP and of the link with the peasantry to an audience some of whose foreign members were plainly sceptical of this interpretation of proletarian revolution. He admitted that the progress of world revolution had not been ``in the straight line which we expected'', and recommended ``a profound study of its concrete development''. Trotsky remarked that, while in 1919 world revolution had seemed ``a question of months'', it was now ``perhaps a question of years''. Practical caution had replaced the unrestrained enthusiasm of the previous congress. 
% 注：去除了 "^ad evaporated", "/maintained", "r \ T v", "V previous" 等乱码和分行符。

Much time was spent analysing the failure of the ``March action'' in Germany, and the internecine divisions in the Italian Left. The ``21 conditions'' of admission to Comintern drawn up by the second congress had split some of the principal foreign parties, and led to the exclusion of sympathizers who were not prepared to accept the rigid discipline imposed by them. Once the first revolutionary wave had receded, only a minority of workers in western countries felt any specific allegiance to communism. The danger was seen that communist parties might degenerate into small sects bound by rigid doctrine and isolated from the main body of workers; the British and American parties, in particular, were warned that it was ``a matter of life and death not to remain a sect''. A new emphasis was placed on the need to woo ``the masses''. Six months after the congress IKKI issued a pronouncement on ``the united workers' front''. This was a call to communists to cooperate with other workers and members of Left parties on joint platforms for specific purposes. Since, however, it was an imperative condition that communists should not sacrifice their independence or their right to criticize, the conception of the united front remained ambiguous, and gave rise to much friction and misunderstanding in the years to come. 
% 注：去除了 "[[/to remain", "I/to woo" 中的符号噪点。

The new turn in foreign policy which accompanied the introduction of NEP extended to Soviet relations with eastern countries. Treaties with Afghanistan and Persia were signed in February 1921, and a treaty with Turkey on the same day as the Anglo-Soviet agreement, March 16, 1921. The Persian treaty seemed difficult to reconcile with the support given at that very moment by Soviet agents to a rebel leader who was seeking to set up an independent republic in northern Persia. But during the summer this support was withdrawn, and the revolt collapsed. The Turkish treaty, which proclaimed the solidarity of the two countries ``in the struggle against imperialism'', caused greater and more lasting embarrassment. Three months before its signature the leader of the illegal Turkish Communist Party had been murdered, and other Turkish communists killed or arrested, by Kemal's agents; and the suppression of communism was a well-advertised aim of Kemal's regime. This was glossed over in the common interest of resistance to British intervention in Turkey. The obligation undertaken in the Anglo-Soviet treaty to refrain from propaganda against the British Empire in Asia also imposed some measure of public restraint. Though Lenin assured the third congress of Comintern that ``the revolutionary movement among hundreds of millions of oppressed peoples of the east grows with remarkable vigour'', the congress itself, unlike its predecessor, was silent on eastern affairs. Lenin, in the peroration of his last speech to a congress of Comintern in November 1922 (he was now a sick man), concluded that ``the most important task in the period now beginning is to study in order to achieve organization, structure, method and content in revolutionary work''. It was a low-key ending. 
% 注：去除了 "(introduction" 和 "tof Kemal's" 等符号错位。

On the other hand, the Soviet Government appeared more decisively than hitherto as the defender of traditional Russian interests. For an almost land-locked country the passage from the Black Sea through the Straits to the Mediterranean had always been a sensitive point. The Soviet-Turkish treaty of March 16, 1921, had guaranteed free passage ``for the commerce of all nations''. But the crux was the passage of warships. Turkey had protested against the use of the Straits by foreign warships without Turkish permission as an infringement of her sovereignty. Soviet Russia, with her depleted naval forces and fear of foreign incursions in the Black Sea, vigorously endorsed the protest. A conference was to meet at Lausanne in the autumn of 1922 to settle terms of peace between the western Powers and Turkey, at which this issue would inevitably be raised; and, rather unexpectedly, the Soviet Government was invited to participate ``in the discussion of the question of the Straits''. Chicherin led the Soviet delegation, and his confrontation with Curzon, then regarded as the main champion of British imperialism in the east, was widely publicized. The issue was settled by a compromise; and the Soviet Government signed, but never ratified, the resulting convention. What had been achieved was general recognition of Soviet Russia as the heir to the rights and interests of the former Russian Empire. 
% 注：去除了 ".warships" 的前导点，并修正 "ler sovereignty" 为 "her sovereignty"。